\documentclass[11pt,a4paper]{article}

% ---------------------------------------------------------
% PREAMBLE (stable, warning-free, Overleaf-safe)
% ---------------------------------------------------------
\usepackage[utf8]{inputenc}
\usepackage[T1]{fontenc}
\usepackage{lmodern}
\usepackage{amsmath, amssymb, amsfonts}
\usepackage{bm}
\usepackage{geometry}
\usepackage{graphicx}
\usepackage{hyperref}
\usepackage{cite}
\usepackage{authblk}
\usepackage{physics}
\usepackage{mathtools}

\geometry{margin=2.4cm}

\title{\textbf{Companion LXV: Fréchet Means of Persistence Diagrams in the Relaxation-Driven Cyclic Cosmology}}
\author{Michael Lehmann \\ \texttt{mi.lehmann@gmx.de}}
\date{April 2026}

\begin{document}
\maketitle

\begin{abstract}
Fréchet means provide a statistical summary of persistence diagrams by minimizing 
the expected Wasserstein distance to an ensemble of diagrams. In weak-lensing 
analyses, they characterize the average topological structure of convergence 
fields across realizations, redshift bins, or smoothing scales. In the standard 
cosmological model, Fréchet means reflect the nonlinear matter distribution and the 
projection of large-scale structure. In the Relaxation-Driven Cyclic Cosmology 
(RDCC), the scale-dependent evolution of the gravitational potential modifies the 
ensemble topology, shifting the mean diagram in a characteristic way. This 
Companion develops a continuous description of Fréchet means in RDCC, deriving how 
the relaxation scale influences ensemble topology, diagram geometry, and the 
observational signatures in weak-lensing surveys. The resulting framework provides 
a distinctive ensemble-level probe of the RDCC gravitational field.
\end{abstract}
% ---------------------------------------------------------
\section{Introduction}

Persistence diagrams encode the birth and death of topological features in 
weak-lensing convergence fields. While Wasserstein distances compare individual 
diagrams, Fréchet means summarize entire ensembles by minimizing the expected 
distance to all diagrams in the set.

In the standard cosmological model, Fréchet means reflect the nonlinear matter 
distribution and the projection of large-scale structure. In RDCC, the 
gravitational potential evolves differently. The relaxation mechanism suppresses 
the decay of small-scale modes and enhances the coherence of large-scale modes. 
This modifies the ensemble topology in a scale-dependent manner, producing a 
distinctive signature in Fréchet-mean analyses.
% ---------------------------------------------------------
\section{Fréchet Means of Persistence Diagrams}

Given an ensemble of persistence diagrams $\{D_{i}\}$, the Fréchet mean $D^{*}$ is 
defined as

\begin{equation}
    D^{*}
    = \arg\min_{D}
      \sum_{i}
      W_{p}(D, D_{i})^{p},
\end{equation}

where $W_{p}$ is the $p$-Wasserstein distance.

The Fréchet mean represents the diagram that best summarizes the ensemble.
% ---------------------------------------------------------
\section{RDCC Modification of Persistence Diagrams}

In RDCC, the convergence evolves as
\begin{equation}
    \kappa_{\mathrm{RDCC}}(\ell)
    = \kappa(\ell)
      \left[
        1 - \chi_{\kappa}
        \left(\frac{\ell}{\ell_{\mathrm{rel}}}\right)^{\alpha}
      \right].
\end{equation}

This modifies the birth and death thresholds of topological features:

\begin{align}
    b_{\mathrm{RDCC}} &= b \left[ 1 - \chi_{b} (\ell/\ell_{\mathrm{rel}})^{\alpha} \right], \\
    d_{\mathrm{RDCC}} &= d \left[ 1 - \chi_{d} (\ell/\ell_{\mathrm{rel}})^{\alpha} \right].
\end{align}

The ensemble of RDCC diagrams becomes
\begin{equation}
    \{ D_{i,\mathrm{RDCC}} \}
    = \left\{
        (b_{i,\mathrm{RDCC}}, d_{i,\mathrm{RDCC}})
      \right\}.
\end{equation}
% ---------------------------------------------------------
\section{Fréchet Means in RDCC}

The RDCC Fréchet mean is

\begin{equation}
    D^{*}_{\mathrm{RDCC}}
    = \arg\min_{D}
      \sum_{i}
      W_{p}(D, D_{i,\mathrm{RDCC}})^{p}.
\end{equation}

To first order in the relaxation parameter,

\begin{equation}
    D^{*}_{\mathrm{RDCC}}
    = D^{*}
      \left[
        1 - \chi_{F}
        \left(\frac{\ell}{\ell_{\mathrm{rel}}}\right)^{\alpha}
      \right].
\end{equation}

This modification reflects the relaxation-driven shift in ensemble topology.
% ---------------------------------------------------------
\section{Ensemble-Level Signatures of RDCC}

The relaxation mechanism enhances the coherence of large-scale modes. Fréchet 
means therefore exhibit:

\begin{itemize}
    \item a shift toward longer-lived large-scale features,
    \item suppression of short-lived small-scale features,
    \item reduced ensemble variance at large scales,
    \item enhanced variance at small scales,
    \item a characteristic turnover near $\ell_{\mathrm{rel}}$,
    \item modified clustering of diagrams in Fréchet space.
\end{itemize}

These signatures provide a direct ensemble-level probe of the RDCC gravitational 
field.
% ---------------------------------------------------------
\section{Conclusion}

Fréchet means of persistence diagrams in the Relaxation-Driven Cyclic Cosmology 
reflect the scale-dependent evolution of the gravitational potential and the 
nonlinear matter distribution. The relaxation mechanism modifies the ensemble 
topology in a scale-dependent manner, producing distinctive signatures in 
weak-lensing surveys. These effects provide a direct observational window into the 
relaxation scale and the temporal structure of the RDCC gravitational field. The 
present Companion establishes the theoretical foundation for future studies of 
topological ensemble statistics, nonlinear structure, and the interplay between 
light and gravity in RDCC.

\bibliographystyle{apsrev4-2}
\nocite{*}
\bibliography{references}


\end{document}
